% Los dos diagramas de casos de uso (UML), uno encima del otro en una página propia y a la misma escala.
% Arriba: configuración del puesto, antes de conducir. El piloto participa en todos salvo en CU6, que solo hace
% el constructor; CU5 lo comparten. Abajo: conducción y análisis. CU10 y CU14 extienden a CU9, y el simulador
% participa en lanzar el juego y en la sesión.
\begin{figure}[p]
	\tikzset{
		cu/.style={draw, thick, ellipse, align=center, minimum width=4.0cm, minimum height=0.72cm, inner sep=0.5pt, font=\tiny, fill=white},
		limite/.style={draw, thick, fill=gray!6, inner sep=4pt},
		titulo/.style={font=\scriptsize\bfseries},
		nombre/.style={font=\scriptsize, align=center},
		dep/.style={-{Stealth[length=1.6mm]}, dashed, thick},
		estereo/.style={font=\tiny, fill=gray!6, inner sep=0.5pt},
		actor/.pic={
			\draw[thick] (0,0.5) circle (0.1);
			\draw[thick] (0,0.4) -- (0,0.12);
			\draw[thick] (-0.2,0.3) -- (0.2,0.3);
			\draw[thick] (0,0.12) -- (-0.15,-0.12);
			\draw[thick] (0,0.12) -- (0.15,-0.12);
		}
	}
	\centering
	\begin{minipage}{\textwidth}
		\centering
		\scalebox{1.25}{\begin{tikzpicture}
			% Mismo ancho de dibujo y de recuadro en las dos figuras, para que queden alineadas
			\path (-0.55,3) (7.25,3);
			\coordinate (izq) at (1.3,3); \coordinate (der) at (5.5,3);
			\node[titulo] (t) at (3.4,6.65) {VH SimGrid};
			\node[cu] (c1)  at (3.4,6.0)  {CU1. Realizar la primera\\puesta en marcha};
			\node[cu] (c2)  at (3.4,5.15) {CU2. Configurar\\el volante};
			\node[cu] (c3)  at (3.4,4.3)  {CU3. Calibrar\\la pedalera};
			\node[cu] (c4)  at (3.4,3.45) {CU4. Configurar las\\luces del aro};
			\node[cu] (c15) at (3.4,2.6)  {CU15. Configurar la\\pantalla del aro};
			\node[cu] (c7)  at (3.4,1.75) {CU7. Gestionar perfiles\\y copias de seguridad};
			\node[cu] (c5)  at (3.4,0.9)  {CU5. Actualizar el firmware\\de un periférico};
			\node[cu] (c6)  at (3.4,0.05) {CU6. Flashear el firmware\\inicial del volante};
			\begin{scope}[on background layer]
				\node[limite, fit=(t)(c1)(c6)(izq)(der)] {};
			\end{scope}

			\pic at (0.2,3.6) {actor};
			\node[nombre, below] at (0.2,3.4) {Piloto};
			\pic at (0.2,0.3) {actor};
			\node[nombre, below] at (0.2,0.1) {Constructor};

			\foreach \c in {c1,c2,c3,c4,c15,c7,c5} \draw (0.45,3.9) -- (\c.west);
			\draw (0.45,0.6) -- (c5.west);
			\draw (0.45,0.6) -- (c6.west);
		\end{tikzpicture}}
		\caption[Diagrama de casos de uso de la configuración del puesto]{Diagrama de casos de uso de la configuración del puesto. (Fuente propia)}
		\label{fig:cu-configuracion}
	\end{minipage}
	\par\vspace{1cm}
	\begin{minipage}{\textwidth}
		\centering
		\scalebox{1.25}{\begin{tikzpicture}
			% Mismo ancho de dibujo y de recuadro en las dos figuras, para que queden alineadas
			\path (-0.55,3) (7.25,3);
			\coordinate (izq) at (1.3,3); \coordinate (der) at (5.5,3);
			\node[titulo] (t) at (3.4,6.65) {VH SimGrid};
			% CU14 encima y CU10 debajo de CU9, para que las dos «extend» sean flechas rectas hacia él
			\node[cu] (c8)  at (3.4,6.0)  {CU8. Lanzar\\un juego};
			\node[cu] (c14) at (3.4,4.95) {CU14. Configurar el equipo\\desde la pantalla del aro};
			\node[cu, minimum width=4.2cm, minimum height=0.95cm] (c9) at (3.4,3.75) {CU9. Conducir una\\sesión asistida};
			\node[cu] (c10) at (3.4,2.55) {CU10. Consultar al\\ingeniero por voz};
			\node[cu] (c11) at (3.4,1.65) {CU11. Planificar la\\estrategia de carrera};
			\node[cu] (c12) at (3.4,0.8)  {CU12. Diseñar\\un dash};
			\node[cu] (c13) at (3.4,-0.05) {CU13. Analizar una\\sesión grabada};

			\draw[dep] (c14) -- node[estereo, right] {\guillemotleft extend\guillemotright} (c9);
			\draw[dep] (c10) -- node[estereo, right] {\guillemotleft extend\guillemotright} (c9);
			\begin{scope}[on background layer]
				\node[limite, fit=(t)(c8)(c9)(c13)(izq)(der)] {};
			\end{scope}

			\pic at (0.35,2.95) {actor};
			\node[nombre, below] at (0.35,2.75) {Piloto};
			\pic at (6.45,4.4) {actor};
			\node[nombre, below] at (6.45,4.2) {Simulador};

			\foreach \c in {c8,c14,c9,c11,c12,c13} \draw (0.6,3.25) -- (\c.west);
			\draw (6.2,4.7) -- (c8.east);
			\draw (6.2,4.7) -- (c9.east);
		\end{tikzpicture}}
		\caption[Diagrama de casos de uso de la conducción y el análisis]{Diagrama de casos de uso de la conducción y el análisis. (Fuente propia)}
		\label{fig:cu-conduccion}
	\end{minipage}
\end{figure}
